\section{问题分析}

\subsection{问题的整体分析}

一份多核方案要回答三件事：每个可分配操作属于哪个子图，每个子图放在哪个核心，同一核心上的子图按什么顺序执行。这三件事互相牵制。子图边界决定哪些张量要经 DDR 中转，分核决定负载是否均衡、哪些依赖变成跨核依赖，同核排序又影响依赖等待和张量在缓存中的驻留时间。切得细，并行机会多，但边界搬运和同步等待随之增加；切得粗，核内复用好，但并行度下降，大子图还会因为缓存放不下而产生换入换出。带通信代价的有向无环图调度已有大量研究\cite{kwok1999static}，本题的特殊之处在于代价由固定的核内调度器和共享带宽模拟共同决定，经典模型不能直接套用。

评价方案的 Makespan 由官方评估程序给出。评估程序先按场景规则插入边界搬运，再用题目附录 C 的启发式确定每个 Task 内的操作顺序，最后在四条流水线和共享带宽池上做离散事件模拟。这个过程没有闭式表达式，也无法对方案求导，但它是确定的：同一方案重复评估结果相同。因此本文的建模分两步。第一步是可以写成数学规划的简化模型，只保留计算量、依赖、边界搬运和固定同步延迟，用来分析代价从哪里来、给出设计原则和性能参照；第二步以官方评估程序作为真实目标，在构造算法生成的候选方案中逐个评估、择优。评估程序在大用例上一次要运行数分钟，候选方案不能靠随机枚举产生，必须由针对结构设计的构造算法给出。

测试用例的结构差别很大，这一点决定了求解框架。有的计算图由几十上百个互不相连的分量组成，按分量分给不同核心就几乎没有跨核通信；有的计算图只有一个连通分量，任何切分都要切断数据依赖，只能按层次切开才能并行；还有的介于两者之间，由几条并行链汇聚而成。不同结构需要不同的切法，参数的合适取值也相差几个数量级，所以本文为不同结构设计不同的构造算法，对每个（用例，核数）组合生成候选方案并由评估程序选出最好的一个，再在其上做保持合法性的改进。三个问题共用同一套决策变量和合法性约束，区别在于代价结构，问题二、三的模型都在问题一的基础上修正得到。

\subsection{问题一的分析}

场景 A 中每个子图是一个独立 Task，Task 之间不保留核内数据。跨子图的张量必须先由生产它的 Task 写回 DDR，再由使用它的 Task 读入，即使两个 Task 在同一核心上相邻执行也是如此。因此问题一的代价有三个来源：一是切图带来的边界搬运，它和其他核心的搬运争用同一份 DDR 带宽；二是 Task 激活等待，同核前一个 Task 结束后至少 100 周期、跨核前驱结束后至少 1000 周期；三是子图过大时核内缓存装不下，调度器插入的换入换出。前两项随切分变细而增加，第三项随切分变粗而增加，负载均衡又要求子图足够多，这些因素需要在同一个模型里权衡。

进一步看，激活等待是沿依赖链逐段累积的：一条依赖链被切成几段、其中几段跨了核，就要付出几次等待。所以真正影响 Makespan 的，不是子图总数，而是依赖路径上被切开了多少次。互不相连的部分切成很多子图并不增加任何一条路径的长度，而同一条长链上反复切分则会让等待线性叠加。建模时需要把边界搬运量和路径上的累积等待写成显式表达式，求解时则要在不拉长依赖路径的前提下换取并行度。题目还要求算法在 2--5 核下都稳定且高效，这意味着既不能靠对每个用例长时间搜索，也要保证核数增加时结果不变差。

\subsection{问题二的分析}

场景 B 把同一核心上的全部子图合成一个 Task。同核子图之间的数据可以留在 L1/UB 中直接使用，只有跨核依赖才需要 COPY\_OUT、500 周期同步和 COPY\_IN。与问题一相比，边界搬运从“子图边界”变为“核边界”，按子图计的激活等待也不再存在，时间代价只与依赖路径上的换核次数有关，同核上切出多少子图本身不付代价。可以预期，场景 B 下合适的切分比问题一更细，问题一的构造算法可以沿用，但参数的合适区域会向细分方向移动。

另一方面，合并成一个大 Task 也带来新的约束。一个张量若在第 $i$ 个子图生产、在第 $j$ 个子图使用，它要跨过中间所有子图一直驻留在缓存中，驻留数据多了，L1/UB 容量就会不够，核内调度器随之插入换入换出。题目要求严格满足 L1/UB 容量约束，所以模型必须加入缓存驻留约束，同核子图的排序也不再只是满足依赖，还直接决定驻留压力。问题二的模型因此需要在问题一的基础上修正通信量和时间约束，并增加驻留约束；求解上除沿用构造算法外，还要重新分析后处理在场景 B 下的作用。

\subsection{问题三的分析}

问题三在场景 B 上增加共享只读 L2 Cache。它只服务 COPY\_IN，未命中的读取完成后才写入缓存，按先进先出淘汰，命中不改变淘汰顺序。于是一次读取能否命中，不仅取决于缓存容量，还取决于同一张量上一次被读取是否已经完成：若多个核心几乎同时读同一张量，第一次读取尚未完成，后面的读取都会落空。模型需要同时刻画命中的容量条件和时序条件，并据此区分哪些缺失是计算图决定的、哪些可以通过调度改善。

命中省下的是 DDR 带宽，只有 DDR 确实繁忙时才能换来时间，所以还需要把缓存收益和 DDR 负载联系起来，估计缓存最多能带来多少加速。题目要求研究容量和带宽的影响，可以固定一批方案、只改变这两个参数重新评估来回答。求解上，问题三与问题二同属场景 B，方案格式和合法性要求相同，问题二的候选方案都可以交给问题三的评估程序复评，再在缓存评估下做后处理，这样既能利用问题二的结果，又能针对缓存的时序特点做改进。
